\section*{Chapter \ref{ch:g11:cost-of-an-algorithm} --- The Cost of an Algorithm}

\begin{solution}{exo:g11:cost-of-an-algorithm:1}
$500$, when the name is absent, or is the last one: every element is compared.
\end{solution}

\begin{solution}{exo:g11:cost-of-an-algorithm:2}
$1000$, $500$, $250$, $125$, $62$, $31$, $15$, $7$, $3$, $1$: nine halvings, $h(1000) = 9$. Check:
$2^9 = 512 \le 1000 < 1024 = 2^{10}$.
\end{solution}

\begin{solution}{exo:g11:cost-of-an-algorithm:3}
$19$ for the sorted list ($n - 1$) and $190$ for the reversed one ($20 \times 19 / 2$); the last column prints $190$
too. The shuffled column depends on the shuffle.
\end{solution}

\begin{solution}{exo:g11:cost-of-an-algorithm:4}
For each \code{i}, the inner body runs \code{i} times: $0 + 1 + \dots + 9 = 45$ for $n = 10$, and
$0 + 1 + \dots + (n - 1) = n(n - 1)/2$ in general, by the pairing of exercise
\ref{exo:g11:cost-of-an-algorithm:6}.
\end{solution}

\begin{solution}{exo:g11:cost-of-an-algorithm:5}
\smallskip\noindent\hspace*{2em}\code{def count_comparisons(sort, t):}

\noindent\hspace*{4em}\code{count = 0}

\noindent\hspace*{4em}\code{def less(a, b):}

\noindent\hspace*{6em}\code{nonlocal count}

\noindent\hspace*{6em}\code{count += 1}

\noindent\hspace*{6em}\code{return -1 if a < b else (1 if b < a else 0)}

\noindent\hspace*{4em}\code{key = functools.cmp_to_key(less)}

\noindent\hspace*{4em}\code{sort([key(v) for v in t])}

\noindent\hspace*{4em}\code{return count}

\smallskip\noindent Each comparison the sort makes between two keys calls \code{less} once. On a sorted list of $10$
elements, \code{sorted} makes $9$.
\end{solution}

\begin{solution}{exo:g11:cost-of-an-algorithm:6}
Turn $i$ compares place $i$ with the $n - 1 - i$ places after it, whatever the values: in all
$S = (n - 1) + (n - 2) + \dots + 1$. Written backwards, $S = 1 + 2 + \dots + (n - 1)$. Adding the two lines term by term
gives $n - 1$ sums, each equal to $n$: $2S = n(n - 1)$, so $S = n(n - 1)/2$.
\end{solution}

\begin{solution}{exo:g11:cost-of-an-algorithm:7}
Printing a list writes every element as text and sends it to the screen, which takes far longer than a comparison. The
timing then measures mostly the printing, which grows with the list printed at each turn, and depends on the terminal,
not on the sort. Time the sort alone, and print afterwards.
\end{solution}

\begin{solution}{exo:g11:cost-of-an-algorithm:8}
Tripling $n$ multiplies a quadratic \omterm{def:g11:cost-of-an-algorithm:cost}{cost} by $3^2 = 9$: about $18$ seconds.
\end{solution}

\begin{solution}{exo:g11:cost-of-an-algorithm:9}
\smallskip\noindent\hspace*{2em}\code{def has_duplicate(t):}

\noindent\hspace*{4em}\code{for i in range(len(t)):}

\noindent\hspace*{6em}\code{for j in range(i + 1, len(t)):}

\noindent\hspace*{8em}\code{if t[i] == t[j]:}

\noindent\hspace*{10em}\code{return True}

\noindent\hspace*{4em}\code{return False}

\smallskip\noindent Without a duplicate, every pair is compared once: $n(n - 1)/2$ comparisons.
\end{solution}

\begin{solution}{exo:g11:cost-of-an-algorithm:10}
A list in increasing order, all different: every \code{t[i]} is larger than the best so far, so the assignment runs at
each of the $n - 1$ turns.
\end{solution}

\begin{solution}{exo:g11:cost-of-an-algorithm:11}
\smallskip\noindent\hspace*{2em}\code{def h(n):}

\noindent\hspace*{4em}\code{count = 0}

\noindent\hspace*{4em}\code{while n > 1:}

\noindent\hspace*{6em}\code{n //= 2}

\noindent\hspace*{6em}\code{count += 1}

\noindent\hspace*{4em}\code{return count}

\smallskip\noindent The check is a loop over \code{range(1, 5001)} with \code{assert 2**k <= n < 2**(k + 1)}; the
tests run it.
\end{solution}

\begin{solution}{exo:g11:cost-of-an-algorithm:12}
Let $w = \mathit{hi} - \mathit{lo}$ be the window's size, $n$ at the start. A turn that does not return leaves either
$\mathit{mid} - \mathit{lo} = \lfloor w/2 \rfloor$ elements, or $\mathit{hi} - \mathit{mid} - 1$, which is
$\lceil w/2 \rceil - 1 \le \lfloor w/2 \rfloor$: at most half of $w$, rounded down. So after $k$ turns the window has at
most as many elements as $n$ halved $k$ times, rounding down each time. After $h(n)$ turns that is $1$, by the
definition of $h(n)$; a window of $1$ element allows one more turn, after which it is empty and the loop stops. In all,
at most $h(n) + 1$ turns, hence at most $h(n) + 1$ elements \code{t[mid]} looked at.
\end{solution}

\begin{solution}{exo:g11:cost-of-an-algorithm:13}
\begin{center}
\begin{tabular}{lrrr}
  \omterm{def:g11:cost-of-an-algorithm:cost}{cost} & $n = 10$ & $n = 1000$ & $n = 10^6$ \\ \hline
  $1000$ & $1000$ & $1000$ & $1000$ \\
  $h(n)$ & $3$ & $9$ & $19$ \\
  $100n$ & $1000$ & $100\,000$ & $10^8$ \\
  $n \cdot h(n)$ & $30$ & $9000$ & $1.9 \times 10^7$ \\
  $n^2$ & $100$ & $10^6$ & $10^{12}$
\end{tabular}
\end{center}
At $n = 10$ the smallest is $h(n)$, then $n \cdot h(n)$, $n^2$, and $1000$ and $100n$ tie. At $10^6$,
$n \cdot h(n)$ is still five times below $100n$, although it grows faster: it passes $100n$ only when $h(n)$ passes
$100$, for $n$ beyond $2^{100}$, a size no list will ever have. How a \omterm{def:g11:cost-of-an-algorithm:cost}{cost} grows tells which wins for large enough $n$;
the \omterm{def:g11:cost-of-an-algorithm:growth}{constants} tell what ``large enough'' means.
\end{solution}

\begin{solution}{exo:g11:cost-of-an-algorithm:14}
\smallskip\noindent\hspace*{2em}\code{def duplicates_by_sorting(t):}

\noindent\hspace*{4em}\code{s = sorted(t)}

\noindent\hspace*{4em}\code{return any(s[k] == s[k + 1] for k in range(len(s) - 1))}

\smallskip\noindent The scan makes at most $n - 1$ comparisons, a linear \omterm{def:g11:cost-of-an-algorithm:cost}{cost}; the sort is the larger part, as the
weekend problem measures.
\end{solution}

\begin{solution}{exo:g11:cost-of-an-algorithm:15}
From $4000$ to $40\,000$ the size is multiplied by $10$, so a quadratic time by $100$: about $100$ times the time printed
for $n = 4000$, a little under $20$ seconds on the author's laptop. The prediction may be off: a list ten times larger
no longer fits in the processor's fastest memory, the machine may be busy, and the small terms of the \omterm{def:g11:cost-of-an-algorithm:cost}{cost}, negligible
at large $n$, were not negligible in the measurement at $4000$.
\end{solution}

\begin{solution}{pb:g11:cost-of-an-algorithm:1}
\textbf{1.} It can answer \code{False} only after checking every pair; with a duplicate it may stop at the first equal
pair.

\textbf{2.} Each pair $i < j$ once: $n(n - 1)/2$. For $n = 100$, $4950$, as printed.

\textbf{3.} Quadratic.

\textbf{4.} It wraps a number and counts, in a variable shared by all its values, each call to \code{<} or \code{==}.
\code{sorted} compares the elements with \code{<}, so each of its comparisons calls \code{Counted.__lt__} once and is
counted.

\textbf{5.} Its count is known exactly, $4\,999\,950\,000$, by the formula; running five billion counted comparisons in
Python would take hours.

\textbf{6.} If \code{s[i] == s[j]} with $i < j$, everything between them in the sorted list lies between the two equal
values, so is equal to them; in particular \code{s[i + 1] == s[i]}.

\textbf{7.} $n - 1$, one per pair of neighbours, when there is no duplicate.

\textbf{8.} At $n = 10$: $32 - 9 = 23$ comparisons for the sort. At $n = 100$: $631 - 99 = 532$.

\textbf{9.} From $n = 7$, for every size up to $100$ on these lists. For small $n$, the sort's own comparisons plus the
$n - 1$ of the scan are more than all the pairs: at $n = 5$, $11$ against $10$.

\textbf{10.} $4\,999\,950\,000 / 1\,628\,843 \approx 3070$ times fewer.

\textbf{11.} $8.0$, $32.7$ and $129.2$~ms: each doubling multiplies the time by about $4$.

\textbf{12.} $100\,000 / 4000 = 25$, and a quadratic time grows like the square: $25^2 = 625$.

\textbf{13.} About $81$~s, a little under a minute and a half.

\textbf{14.} Building the set reads each of the $n$ numbers once: its \omterm{def:g11:cost-of-an-algorithm:cost}{cost} is linear, it grows with $n$.

\textbf{15.} The set, $5.1$~ms; sort then scan, $20.2$~ms; every pair, about $81$~s, predicted.

\textbf{16.} $81 / 0.0202 \approx 4000$, as printed.

\textbf{17.} A comparison \omterm{def:g11:cost-of-an-algorithm:cost}{costs} a different time in each: \code{sorted} compares numbers inside Python's own fast
machinery, while \code{by_pairs} runs two Python loops around each comparison; the sort also moves elements, which is
counted in time but not in comparisons. Counts predict how a time grows, not the ratio between two programs.

\textbf{18.} For a thousand numbers checked once, $8$~ms against a fraction of a millisecond: it hardly matters. For a
million, the pairs would take $250^2 = 62\,500$ times the time at $n = 4000$, more than two hours, against a fraction
of a second.

\textbf{19.} The time also has smaller terms that the square neglects, and a large list behaves differently in the
machine's memory; the prediction extends a law measured from $1000$ to $4000$ far beyond where it was measured.

\textbf{20.} On these shuffled lists of different numbers, \textbf{sort then scan needs fewer comparisons than every
pair from $n = 7$ on}; at $n = 100\,000$ it was measured \textbf{about $4000$ times faster} on the author's laptop
($20.2$~ms against about $81$~s predicted for every pair), and a set was faster still.
\end{solution}
